\section*{Executive Summary}
\addcontentsline{toc}{section}{Executive Summary}
\label{sec:executive}

Olist is a Brazilian marketplace that places small sellers on large retail
platforms and manages fulfilment for them; its public dataset --- 99{,}441
orders in nine relational tables, 2016--2018 --- records each order from
checkout to delivery to review. This project asks two questions of it, at two
moments in an order's life, and turns the answers into decisions.
\textbf{Problem 1, at checkout:} how long will this order take, and will it miss
the date the customer was promised? \textbf{Problem 2, at delivery:} will this
parcel draw a one- or two-star review? The first serves the operations team,
which can re-quote a promise, warn a buyer or expedite a parcel; the second
serves customer experience, which can reach a customer before the review is
written. The two are halves of one chain: a late delivery raises the
probability of a low review from 9.2\% to 54.0\%, the strongest relationship in
the data.

Three disciplines run through the work. \emph{Leakage is defined per prediction
moment}, and a guard enforces the right feature set before any model is fitted
--- one forbidden column lifts the late flag's ROC-AUC from 0.589 to 1.000.
\emph{Time is respected}: the split is chronological and customer-disjoint,
cross-validation uses expanding windows, and the test window (July--August
2018) was scored once, for models already chosen. \emph{Complexity has to earn
its place}: three model families --- linear, tree-based, ensemble --- are run
as ladders from their simplest variant up, and every rung is compared with the
one below it.

The lead-time regression reaches a mean absolute error of \textbf{3.90 days},
41\% below a constant fitted on the past and 60\% below the platform's own
promised date. Thresholding that forecast against the promise gives a late-order
flag at 2.9 times the no-skill floor, where the best classifier trained on the
binary label reaches 1.2 times. The low-review model reaches ROC-AUC 0.749 and
PR-AUC 0.414, \textbf{3.75 times its floor}, on a target where always predicting
``satisfied'' scores 89\% accuracy and catches nothing. Read as a ranked list,
which is how a team with limited capacity uses it, contacting the riskiest 5\%
of deliveries reaches a would-be low reviewer two times in three.

Three recommendations follow, each with a stated cost. Operations should
re-quote every promise with the model's estimate and pre-warn buyers whose late
risk exceeds 11\%, both nearly free; it should \emph{not} expedite on the flag,
whose precision never clears the 0.30 a R\$12 intervention would need. Customer
experience should work each day's deliveries from the top of the review-risk
list to its capacity: at illustrative costs of R\$5 per contact and R\$40 per
low review the policy is worth about R\$1{,}466 per thousand delivered orders.
Because that figure assumes a contact prevents the review it reaches ---
something no historical data can show --- the report sizes the pilot that would
test it: a randomised hold-out inside the daily list, about six weeks at ten
contacts a day, with break-even at roughly one averted review in five.

Everything is deployed and reproducible. A Streamlit application at
\url{https://team10-it5006-olist.streamlit.app} serves the three models on
single orders and batch files, returns the ranked triage list, replays any day
of the held-out test window with the real outcomes revealable, and carries a
model card stating exactly which estimators are deployed. The repository holds
the source tables, ten executed notebooks, the fitted artifacts and the scripts
that regenerate every number in this document; a suite of \checkcount{}
validation checks ties the metadata, the pickled models, the result tables and
this text to one another. Both rounds of evaluator feedback are incorporated
throughout and answered point by point in Appendix~\ref{app:feedback}.
